\section{Discusión}
\label{sec:discusion}

\subsection{La cobertura baja como hallazgo, no como defecto}

El resultado más llamativo es también el más fácil de malinterpretar: \pctIndeterminado{}\%
de los mensajes no recibe etiqueta de polaridad. Una lectura ingenua lo tomar\'\i{}a por
un fallo del m\'etodo. Sostenemos lo contrario. El corpus es, antes que nada, un canal
operativo: sus mensajes sirven para coordinar, informar y resolver dudas. La funci\'on
primordial del lenguaje en ese contexto es referencial, no expresiva. Un m\'etodo que
forzara una etiqueta en cada mensaje estar\'\i{}a, en el mejor de los casos, adivinando, y
en el peor, imputando neutralidad a enunciados que simplemente no hablan de afecto.

La distinci\'on entre \emph{indeterminado} y \emph{neutro} es, por tanto, sustantiva.
Neutro describe un mensaje con se\~nal afectiva cuyo balance es cero (por ejemplo,
``\emph{me gusta pero me preocupa}''); indeterminado describe la ausencia de se\~nal. La
decisi\'on de reservar la etiqueta neutra para el primer caso evita el sesgo m\'as com\'un
en los estudios de sentimiento aplicado: la sobreestimaci\'on de la neutralidad por
defecto. La contrapartida es una cobertura modesta; la asumimos como el precio de una
medici\'on honesta.

\subsection{Un clima predominantemente positivo}

Entre los mensajes que s\'\i{} aportan evidencia, la proporci\'on de positivos
(\pctPositivos{}\%) duplica con creces la de negativos (\pctNegativos{}\%). Este hallazgo
es consistente con la funci\'on social del grupo: un espacio de ayuda mutua donde la
gratitud y el reconocimiento son moneda corriente. La emoci\'on dominante,
\emph{\emocionTop{}}, y la presencia destacada de reacciones como \reactTop{} refuerzan
esa interpretaci\'on. Cabe se\~nalar, sin embargo, un posible sesgo de deseabilidad: los
mensajes con carga negativa fuerte (quejas, conflictos) pueden haberse canalizado por
v\'\i{}as privadas y no aparecer en el registro p\'ublico del grupo.

\subsection{Temas y estructura latente}

El modelado de tópicos confirma la vocaci\'on administrativa de la comunidad. Los temas
recuperados ---pagos y matr\'\i{}culas, horarios y paralelos, inscripciones, comunicados y
consultas--- se corresponden con los hitos del ciclo acad\'emico. Es notable que los
términos asociados a trámites aparezcan pr\'oximos a vocabulario de urgencia, lo que
explica por qué la ansiedad y la frustraci\'on, aun siendo minoritarias, se concentran en
esos temas. La interpretaci\'on debe ser prudente: tal como advierte
\citet{grimmer2013text}, los tópicos son artefactos estad\'isticos y su etiquetado es una
construcci\'on del analista.

\subsection{Dinámica de estados}

Las cadenas de Markov a\~naden una dimensi\'on temporal a la fotograf\'\i{}a est\'atica de
la polaridad. El estado neutro funciona como un atractor: con una permanencia media de
\permanenciaNeutro{} turnos y \pctNeutro{}\% de la masa estacionaria, la conversaci\'on
regresa una y otra vez a su régimen operativo. Los episodios negativos, en cambio, son
breves (recuperaci\'on de \recuperacionNegativa{} turnos). Este patr\'on sugiere un
sistema autorregulado: las irrupciones emocionales no se cronifican, sino que se disuelven
en el flujo cotidiano de la coordinaci\'on acad\'emica.

\subsection{Implicaciones pr\'acticas}

Los resultados tienen aplicaciones concretas. Un coordinador acad\'emico podr\'\i{}a
monitorizar la proporci\'on de mensajes negativos como indicador temprano de fricciones,
prestando atenci\'on a los picos coincidentes con per\'\i{}odos de inscripci\'on o pago.
La identificaci\'on de los temas m\'as asociados a emociones negativas permitir\'\i{}a
anticipar dudas y mejorar la comunicaci\'on institucional. En t\'erminos m\'as generales,
el estudio demuestra que es posible obtener inteligencia cualitativa de un canal
informal sin incurrir en la vigilancia individual: toda la anal\'\i{}tica se realiza sobre
agregados an\'onimos.

\subsection{Comparaci\'on con la literatura}

El balance positivo observado es coherente con los hallazgos de
\citet{golder2011diurnal} y \citet{bollen2011twitter}, quienes documentan una tendencia
general hacia la positividad en redes sociales, con modulaci\'on por hora y día. La
distribuci\'on horaria concentrada en la franja laboral reproduce el patr\'on diurno
descrito por esos trabajos. La contribuci\'on espec\'\i{}fica de este estudio es el
\'enfasis en la separaci\'on de se\~nales y en la honestidad de la cobertura, aspectos
menos atendidos en la literatura orientada a maximizar \emph{accuracy} sobre corpus
etiquetados.
